\documentclass{article}
\usepackage{amsmath,amssymb}

\begin{document}

\section*{Equations from Elkins, Sood \& Rumpf (AAS 20-475)}

\begin{equation}
R_t(s_t) = \sum_{k=0}^{T} \gamma^{t} r_{t+k}
% Paper prints \gamma^{t}; the standard form is \gamma^{k}
\end{equation}

\begin{equation}
Q^{\pi} = \hat{\mathbb{E}}\left[ R_t \mid s_t, a_t \right]
\end{equation}

\begin{equation}
\nabla_{\theta} J(\theta) = \hat{\mathbb{E}}\left[ \nabla_{a_t} Q^{\pi}(s_t, a_t)\, \nabla_{\theta} \pi_{\theta}(s) \right]
\end{equation}

\begin{equation}
Q^{\pi}(s_t, a_t) = r_t + \gamma\, \hat{\mathbb{E}}\left[ Q^{\pi}(s_{t+1}, a_{t+1}) \right]
\end{equation}

\begin{equation}
y = r_t + \gamma\, Q_{\psi'}(s_{t+1}, a_{t+1})
\end{equation}

\begin{equation}
\psi' \longleftarrow \tau \psi + (1 - \tau)\psi'
\end{equation}

\begin{equation}
\mathbf{I} \equiv \operatorname{diag}(I_1, I_2, I_3) = \operatorname{diag}(0.872,\, 0.115,\, 0.797)\ \text{kg m}^2
\end{equation}

\begin{equation}
\vec{M} = \mathbf{I}\dot{\vec{\omega}} + \boldsymbol{\omega}^{\times}\mathbf{I}\vec{\omega}
\end{equation}

\begin{equation}
\boldsymbol{\omega}^{\times} =
\begin{bmatrix}
0 & -\omega_3 & \omega_2 \\
\omega_3 & 0 & -\omega_1 \\
-\omega_2 & \omega_1 & 0
\end{bmatrix}
\end{equation}

\begin{equation}
\mathbf{q}_e =
\begin{bmatrix} \hat{e}\sin(\phi/2) \\ \cos(\phi/2) \end{bmatrix}
=
\begin{bmatrix} \vec{q} \\ q_s \end{bmatrix}
\end{equation}

\begin{equation}
\dot{\mathbf{q}}_e = \frac{1}{2}\boldsymbol{\Omega}(\vec{\omega})\,\mathbf{q}_e
\end{equation}

\begin{equation}
\boldsymbol{\Omega}(\vec{\omega}) =
\begin{bmatrix}
-\boldsymbol{\omega}^{\times} & \vec{\omega} \\
-\vec{\omega}^{T} & 0
\end{bmatrix}
\end{equation}

\begin{equation}
s_t = \left\{ \mathbf{q}_e,\ \dot{\mathbf{q}}_e,\ \vec{\omega} \right\}
\end{equation}

\begin{equation}
r_t =
\begin{cases}
0.1 & q_{s,t} > q_{s,t-1} \\
-0.1 & \text{otherwise}
\end{cases}
\end{equation}

\begin{equation}
r_t = \boldsymbol{\alpha} \cdot \mathbf{q}_e =
\begin{bmatrix} -1 \\ -1 \\ -1 \\ 1 \end{bmatrix}
\cdot
\begin{bmatrix} q_1^2 \\ q_2^2 \\ q_3^2 \\ q_s^2 \end{bmatrix}
\end{equation}

\begin{equation}
\vec{\tau}_{dist} = 0.05
\begin{bmatrix} I_3 - I_2 \\ I_1 - I_3 \\ I_2 - I_1 \end{bmatrix}
=
\begin{bmatrix} 0.0341 \\ 0.0038 \\ -0.0379 \end{bmatrix}\ \text{Nm}
\end{equation}

\section*{Related (unnumbered in the paper)}

TD3 clipped double-Q target:
\begin{equation*}
y = r_t + \gamma \min_{i=1,2} Q_{\psi_i'}\!\left(s_{t+1},\, \pi_{\theta'}(s_{t+1}) + \epsilon\right),
\quad \epsilon \sim \operatorname{clip}(\mathcal{N}(0,\sigma), -c, c)
\end{equation*}

Eq.~15 simplified using the unit-norm constraint:
\begin{equation*}
r_t = q_s^2 - (q_1^2 + q_2^2 + q_3^2) = 2q_s^2 - 1 = \cos\phi
\end{equation*}

Reward-switch threshold vs.\ angle:
\begin{equation*}
q_s \ge 0.999962 \;\Rightarrow\; \phi \le 2\arccos(0.999962) \approx 1.0^\circ
\end{equation*}

\end{document}


// this is from  https://www.researchgate.net/publication/343834157_Adaptive_Continuous_Control_of_Spacecraft_Attitude_Using_Deep_Reinforcement_Learning